\section{论文的数学设置：从材料到测量}
前面的部分给出了背景和直觉。下面按论文的论证顺序展开中文推导：先固定材料和测量的含义，再消去电流坐标，随后推导材料步变化与实际收益。这里整理的是英文正文和理论说明中已有的结果，不增加实验或改变成立条件。

\subsection{物理材料空间与多照明}
用 $\chi$ 表示介电对比度（dielectric contrast），$X$ 表示对应的材料响应算子（material-response operator）。第 $p$ 次照明的感应电流（induced current）满足
\begin{equation}
 j_p=X(e_p^{\rm inc}+G_Dj_p),\qquad L=I-XG_D.
\end{equation}
$G_D$ 是内部矢量 Green 算子（internal vector Green operator），$e_p^{\rm tot}=e_p^{\rm inc}+G_Dj_p$ 是总场（total field）。材料变化用实坐标 $s\in\mathbb R^d$ 表示；复介电材料的实部和虚部都是这个实向量的组成部分。物理材料切向（physical material tangent）$Q_T$ 满足
\begin{equation}
 \delta\chi=Q_Ts,\qquad Q_T^\top M_\chi Q_T=I.
\end{equation}
这里的等式在实化坐标中理解，$M_\chi$ 是体积加权的材料 $L^2$ 内积。这样，同一材料子空间仅仅换单位或旋转坐标，不会人为改变某个方向的重要性。

对电流方程求方向导数（directional derivative），得到
\begin{equation}
 L\,\delta j_p=DX[Q_Ts]e_p^{\rm tot}=:B_ps,
 \qquad K_p=L^{-1}B_p,\qquad J_p=SK_p.
\end{equation}
$B_p$ 是材料注入映射（material-injection map）：它将材料扰动转成电流方程的右端。实际电流变化是 $K_ps$，并不是 $B_ps$。$S$ 包括接收、归一化与需要时的噪声白化（noise whitening）；$J_p$ 是材料到测量的 Jacobian。

所有照明共享同一个材料向量 $s$，但有各自的电流。把各次测量的实部、虚部纵向堆叠（stacking），形成实线性映射 $J$。后面的 $J^*$ 是该实材料内积下的伴随（adjoint）。例如，复输出矩阵 $F$ 作用于实材料坐标时，伴随为 $v\mapsto\operatorname{Re}(F^Hv)$，不能遗漏取实部。电流分块公式可先用复记号写出，再对照明做块对角提升和实化（realification）。

\subsection{DDA 中极化率的导数必须保留}
实际离散偶极子近似（discrete-dipole approximation, DDA）使用极化率（polarizability）$a(\chi)$，而不是把 $X$ 直接替换为 $\chi$。对体积为 $v$ 的单元，所用模型为
\begin{equation}
 a(\chi)=\frac{3v\chi}{\chi+3-i3c_kv\chi},\qquad
 a'(\chi)=\frac{9v}{(\chi+3-i3c_kv\chi)^2},\qquad c_k=\frac{k^3}{6\pi}.
\end{equation}
在偶极矩坐标下，$L=I-D_aG_{\rm off}$，扰动激励包含 $a'(\chi)\delta\chi\,e_p^{\rm tot}$；电流范数再按单元体积归一化。因而材料注入同时依赖当前材料、局部总场和允许的材料方向。把 $a$ 当成对比度，或者漏掉 $a'$，都会改变实际研究的材料导数。完整离散Green核及坐标因子见英文补充材料第1节。

\subsection{冻结状态下的局部材料问题}
固定当前材料、照明、残差和正则化。令 $r$ 为预测减观测的测量残差（data residual），$\Lambda\succ0$ 为正则化及阻尼曲率，$\ell$ 为共同的先验线性项。Gauss--Newton局部二次模型写成
\begin{equation}
 \begin{aligned}
 \Phi(s)&=\tfrac12\|r+Js\|^2
             +\tfrac12s^*\Lambda s+\ell^*s\\
 &=\tfrac12s^*Hs+h^*s+\tfrac12\|r\|^2,\qquad
 H=J^*J+\Lambda,\quad h=J^*r+\ell.
 \end{aligned}
\end{equation}
这里 $H$ 是正则化 Gauss--Newton Hessian，也就是材料法方程（material normal equation）的矩阵，不是包含全部二阶导数的精确非线性 Hessian。完整局部解为 $s_*=-H^{-1}h$，并且
\begin{equation}
 \Phi(s)-\Phi(s_*)=\tfrac12\|s-s_*\|_H^2,
 \qquad \|u\|_H^2=u^*Hu.
\end{equation}
因此，局部二次误差和 $H$-范数步误差等价；它们都不等于最终材料重建误差。

\section{消去保留电流：反馈修饰的导数变化}
\subsection{从分块电流方程得到 Schur complement}
设 $V$ 为原先保留的正交电流基，$C$ 为指定新增候选基（candidate bank），满足 $V^*V=C^*C=I$、$V^*C=0$。令 $A=I-L$，并用 $A_{VC}=V^*AC$ 等记号表示分块。假设 $D_V=I-A_{VV}=V^*LV$ 可逆，原电流响应为
\begin{equation}
 R_0=VD_V^{-1}V^*,\quad K_0=R_0B,\quad J_0=SK_0.
\end{equation}
注意 $R_0$ 作用在电流空间；后面 $M_0=(J_0^*J_0+\Lambda)^{-1}$ 作用在材料空间。

对任意材料右端，增广后的电流写成 $V\alpha+C\gamma$。Galerkin投影要求方程残差与整个保留空间正交，因此
\begin{equation}
 \begin{bmatrix}D_V&-A_{VC}\\-A_{CV}&I-A_{CC}\end{bmatrix}
 \begin{bmatrix}\alpha\\\gamma\end{bmatrix}
 =\begin{bmatrix}B_Vs\\B_Cs\end{bmatrix},
 \qquad B_V=V^*B,\quad B_C=C^*B.
\end{equation}
第一行给出 $\alpha=D_V^{-1}(B_Vs+A_{VC}\gamma)$。代入第二行后，新增坐标满足
\begin{equation}
 \Gamma_C\gamma=N_Cs,
 \quad\Gamma_C=I-A_{CC}-A_{CV}D_V^{-1}A_{VC},
 \quad N_C=B_C+A_{CV}D_V^{-1}B_V.
\end{equation}
$\Gamma_C$ 是 Schur补矩阵（Schur complement），描述新增空间自身作用和新增—保留—新增反馈回路。$N_C$ 则是消去原坐标后，材料真正进入新增方程的激励。

\subsection{反馈修饰切向分解及其物理含义}
定义
\begin{equation}
 \bar C=C+VD_V^{-1}A_{VC}=(I-R_0L)C,\qquad Y_C=S\bar C.
\end{equation}
只要 $\Gamma_C$ 可逆，回代立即得到
\begin{equation}\label{zh:tangent}
 \boxed{K_C-K_0=\bar C\Gamma_C^{-1}N_C,\qquad
 J_C-J_0=Y_C\Gamma_C^{-1}N_C.}
\end{equation}
这就是反馈修饰切向分解（feedback-dressed tangent factorization）。在对比源表述下，$A=XG_D$，于是
\begin{equation}
 \bar C=C+R_0XG_DC,\qquad Y_C=SC+SR_0XG_DC.
\end{equation}
$SC$ 是直接辐射；$G_DC$ 是新增电流产生的内部场；$XG_DC$ 是该场激发的材料电流作用；$R_0$ 进一步处理保留空间中的相互散射。接收器看到两条路径的复振幅叠加，既可能增强，也可能抵消。另一侧 $N_C=B_C+A_{CV}D_V^{-1}B_V$ 表明，材料既能直接激发新增方向，也能先激发保留方向再传到新增方向。

这不是给直接输出乘一个经验增益。一般 $\bar C$ 不再与 $V$ 欧氏正交，却满足 $V^*L\bar C=0$：它正是保持原Galerkin方程不被破坏的修正。整个解释是频域反馈（frequency-domain feedback），不额外假定控制系统的时域稳定性。

\subsection{弱反馈的绝对界与相对界不同}
若 $\|A_{VV}\|_2\le\eta<1$，Neumann级数给出 $\|D_V^{-1}\|_2\le(1-\eta)^{-1}$，从而
\begin{equation}
 \|Y_C-SC\|_2\le\frac{\|SV\|_2\|A_{VC}\|_2}{1-\eta}.
\end{equation}
要相对直接输出给出统一保证，还需要 $SC$ 在所研究系数空间中有正的最小奇异值：
\begin{equation}
 \frac{\|(Y_C-SC)u\|_2}{\|SCu\|_2}
 \le\frac{\|SV\|_2\|A_{VC}\|_2}
 {(1-\eta)\sigma_{\min}(SC)},\qquad u\ne0.
\end{equation}
对于严格位于接收零空间的电流（receiver-dark current），这个正下界不存在；近乎暗的方向即使有很小的正奇异值，相对界也可能大到无法提供有用保证。它的绝对反馈可以不大，却远大于直接输出。因此，大的反馈/直接比值本身不能证明谐振；高材料对比度也不能单独证明反馈一定强。

\section{材料步为什么需要两组方向：闭包定理的中文证明}
令 $T_C=\Gamma_C^{-1}N_C$，于是 $J_C=J_0+Y_CT_C$。在相同 $r,\Lambda,\ell$ 下，定义
\begin{equation}
 \begin{aligned}
 H_0&=J_0^*J_0+\Lambda,& M_0&=H_0^{-1},\\
 s_0&=-M_0(J_0^*r+\ell),&
 s_C&=-(J_C^*J_C+\Lambda)^{-1}(J_C^*r+\ell).
 \end{aligned}
\end{equation}
\paragraph{定理：双侧材料步闭包（paired material-step closure）。}
这里的观测、Schur和注入因子需要一起实化。$P$为照明次数；先对照明作块对角提升 $\mathbf Y=I_P\otimes Y_C$、$\Gamma_{\rm blk}=I_P\otimes\Gamma_C$，纵向堆叠注入 $\mathbf N$，再定义 $\mathcal R(F)=[\operatorname{Re}F,-\operatorname{Im}F;\operatorname{Im}F,\operatorname{Re}F]$ 和 $\widehat N=[\operatorname{Re}\mathbf N;\operatorname{Im}\mathbf N]$。于是 $\widehat D=\mathcal R(\mathbf Y)\mathcal R(\Gamma_{\rm blk})^{-1}\widehat N$，与数据的固定堆叠顺序只差一个排列。以下材料空间公式中的 $Y_C,\Gamma_C,N_C$ 指这组实化因子，伴随为实转置。这一步保留跨照明共享材料，不能把一个复电流方向误当作一个实材料秩。
在上述可逆性、共同正则化和精确算子作用条件下，若没有截掉非零作用方向，则
\begin{equation}\label{zh:closure}
 \begin{aligned}
 s_C-s_0={}&-M_0J_0^*Y_C(T_Cs_C)\\
 &-M_0N_C^*\Gamma_C^{-*}Y_C^*(r+J_Cs_C)
 \in\mathcal E_0(C),\\
 \mathcal E_0(C)&=\operatorname{Range}M_0[J_0^*Y_C,N_C^*].
 \end{aligned}
\end{equation}
\paragraph{证明。}
展开曲率差和线性项差：
\begin{equation}
 \begin{aligned}
 \Delta H&=J_0^*Y_CT_C+T_C^*Y_C^*J_0+T_C^*Y_C^*Y_CT_C,\\
 \Delta h&=T_C^*Y_C^*r.
 \end{aligned}
\end{equation}
两套法方程相减，得到
\begin{equation}
 H_0(s_C-s_0)=-\Delta Hs_C-\Delta h
 =-J_0^*Y_CT_Cs_C-T_C^*Y_C^*(r+J_Cs_C).
\end{equation}
左乘 $M_0$，再使用 $T_C^*=N_C^*\Gamma_C^{-*}$，即得式\eqref{zh:closure}。\hfill$\square$

第一组 $M_0J_0^*Y_C$ 是观测侧材料拉回（observation-side material pullback）：新增观测响应与原材料敏感度发生耦合，改变原有材料方向上的步长。第二组 $M_0N_C^*$ 是注入侧材料拉回（injection-side material pullback）：更新后的测量残差沿新增激励路径返回材料空间。只看右端改变 $\Delta h$ 会漏掉第一组中的曲率耦合。

闭包告诉我们步差可能出现在哪里，并不保证当前步一定改变。若 $r=0,\ell=0$，即使 $J_C-J_0\ne0$，仍有 $s_C=s_0=0$。物理通路给出影响的可能性，当前残差和正则化决定它是否实际参与更新。

\subsection{非单调曲率的最小例子}
取
\begin{equation}
 J_0=[1,0],\quad J_C=[1,1],\quad\Lambda=I,\quad r=1,\quad\ell=0.
\end{equation}
直接求解可得 $s_0=(-1/2,0)^\top$、$s_C=(-1/3,-1/3)^\top$，步差为 $(1/6,-1/3)^\top$。观测侧沿第一个材料坐标，注入侧沿第二个材料坐标，单独一侧都无法表达全部步差。同时
\begin{equation}
 \Delta H=\begin{bmatrix}0&1\\1&1\end{bmatrix},\qquad\det\Delta H=-1.
\end{equation}
因此曲率差是不定矩阵（indefinite matrix）。这个例子可嵌入 $L=I$ 的Galerkin模型，说明不定性不要求多重散射才能发生。Maxwell反馈决定实际物理因子；最小二乘法方程本身产生两侧耦合。增加电流空间是改进已有测量模型，不是增加独立测量，二者不能混为一谈。

\subsection{结构闭包与当前误差对齐是两件事}
设 $Z$ 是 $\mathcal E_0(C)$ 的独立基。在不改变完整局部目标的前提下，令修正步为 $s=s_0+Zq$。记完整目标在 $s_0$ 的梯度缺陷为 $d_0=Hs_0+h$，则
\begin{equation}
 q_*=-\bigl(Z^*HZ\bigr)^{-1}Z^*d_0,
 \qquad
 \Phi(s_0)-\Phi(s_0+Zq_*)
 =\tfrac12 d_0^*Z(Z^*HZ)^{-1}Z^*d_0.
\end{equation}
因为 $s_C$ 在这个仿射空间中，最优修正一定不差于指定的 $s_C$。但增益取决于 $Z^*d_0$；若它为零，修正没有收益。用 $P_Z$ 表示到 $H^{1/2}Z$ 的正交投影，还可把增益写为 $\tfrac12\|P_ZH^{-1/2}d_0\|^2$。这解释了为什么完整描述某个模型变化的方向，不一定是消除当前完整误差的最佳方向。材料Krylov方向直接从误差出发，完全可能更有效。
